\documentclass[11pt,a4paper]{article}
\usepackage[margin=2.4cm]{geometry}
\usepackage[T1]{fontenc}
\usepackage{amsmath,amssymb}
\usepackage{microtype}
\usepackage{xcolor}

\usepackage{parskip}
\usepackage[colorlinks=true,linkcolor=blue!60!black,urlcolor=blue!60!black]{hyperref}

\newcommand{\sowhat}[1]{\emph{So what?} #1}

\title{\Large\textbf{A plain-language summary}\\[6pt]
\large of ``Rule reconstruction from a single spacetime diagram yields
simulation-limited prediction of the finite-horizon damage response of
cellular automata''}
\author{Manuscript authors: Michiel Rollier and Jan M.\ Baetens\\
\normalsize Department of Data Analysis and Mathematical Modelling, Ghent University}
\date{Summary prepared 21 August 2026}

\begin{document}
\maketitle

\noindent\fbox{\parbox{0.965\textwidth}{\textbf{Which version this
summarizes.} This document describes the manuscript as it stands on
21~August~2026: the revision responding to the fifth referee report
(registered protocol revisions 2--14; git branch
\texttt{review-round7-fixes}, commit \texttt{8bf5b4a}; 16-page manuscript
prepared for \emph{Chaos}). Every number quoted below appears verbatim in
that version. If the manuscript is revised again, this summary describes the
older text.}}

\section*{1. The whole paper in one breath}

We ask whether a neural network that looks at a single picture of a simple
simulated universe can predict how that universe reacts to a tiny
disturbance---and we show that, whenever the picture is clean, an
old-fashioned and fully transparent alternative beats it: read the
universe's rule straight off the picture, then re-run the disturbance
experiment yourself. The learned models only start earning their keep once
the picture is corrupted by noise or missing data, and the paper maps
exactly where that happens.

\section*{2. Why should anyone care?}

There is a growing scientific practice of training machine-learning models
to read dynamical properties---chaos, stability, sensitivity to
perturbation---directly off a single observation of a system, skipping the
experiments that define those properties. Published comparisons in that
literature often lack strong classical baselines. This paper builds the
strongest possible baseline \emph{exactly}, in the cleanest laboratory
available, works out what ``as good as possible'' even means when the
correct answers are themselves noisy measurements, and then shows that the
learned shortcut falls far short of that baseline while looking deceptively
good on the usual score.

\sowhat{Before trusting a learned shortcut for predicting how a system
responds to perturbation, check whether the observation already contains the
system's mechanism. If it does, reading the mechanism out and simulating is
the baseline to beat---and in this study, across every configuration tested,
nothing beat it.}

\section*{3. The background you need (five minutes)}

\textbf{Cellular automaton (CA).} A row of 127 cells arranged in a ring,
each cell either black or white. A fixed local \emph{rule} decides each
cell's next colour from its current neighbourhood. The paper uses two rule
families: \emph{elementary} CA, where a cell sees itself and one neighbour
on each side (256 possible rules, of which 88 are genuinely different once
mirror-image and colour-swapped versions are identified---each such family
is called an \emph{orbit}), and \emph{radius-two} CA, where a cell sees two
neighbours on each side (about four billion possible rules, from which 800
are sampled). Despite their simplicity, these systems produce everything
from frozen order to full chaos to intricate interacting structures.

\textbf{Spacetime diagram.} Stack 127 successive rows on top of each other
and you get a $127\times127$ black-and-white image: the complete visual
history of one run. This single image is the only input every predictor in
the paper is allowed to see.

\textbf{Rule versus behaviour.} The field distinguishes the
\emph{genotype}---the local update rule, a small lookup table---from the
\emph{phenotype}---the visible behaviour in the diagram. A key subtlety: a
clean diagram contains, hidden in its pixel-to-pixel transitions, the entire
rule table. Any method meant to judge \emph{behaviour} can secretly succeed
by reading the \emph{rule} instead.

\textbf{Damage spreading (the ``butterfly-effect'' probe).} Run the same CA
twice from the same starting row, but flip a single cell in one copy, and
watch how the difference between the two runs evolves. Four numbers
summarize the outcome: \emph{damage survival} (does any difference remain at
the end?), \emph{damage fraction} (how much of the ring differs?),
\emph{spreading rate} (how fast does the damaged region grow?), and
\emph{cone fill} (how densely damaged is that region?). Each number is an
average over 256 repetitions with random starting rows---so each ``correct
answer'' is itself a noisy measurement, like a lab value with error bars.

\textbf{The task.} These four statistics are defined by \emph{twin} runs, so
no formula computes them from one image. The machine-learning proposition
(``amortization'') is: pay once to train a model, then predict the four
numbers from a single image, instantly, forever after.

\textbf{The scores.} $R^2$ is the standard accuracy score: $1$ is perfect,
$0$ means no better than always guessing the average, negative means worse
than that. Because the correct answers are noisy, even a perfect method
cannot reach $1$---an issue the paper takes very seriously. It therefore
also reports $\rho$: the prediction error measured in units of the answer's
own measurement noise. On that scale, $\rho=\sqrt2\approx1.41$ means ``as
accurate as re-measuring the answer with a fresh simulation'' and $\rho=1$
means ``perfect knowledge of the ideal noise-free answer''.

\section*{4. A walk through the paper}

\subsection*{Abstract, opening quotation, and Introduction (\S I)}

\textbf{What it does.} Sets up the question and announces the three results:
(1) a single clean diagram identifies the exact rule for $100\%$ of
elementary and $97.5\%$ of radius-two test rules, and re-simulating from the
recovered rule is then \emph{provably} as good as re-measuring the answer;
(2) two trained direct predictors fall $3$--$35$ noise units short of that;
(3) the mechanistic route is brittle---modest noise, heavy masking, or
extreme starting densities break it, and the paper maps what to use beyond.

\sowhat{The intellectual move is taking an ``obvious'' trick seriously.
Reading a CA rule out of its diagram is thirty-five-year-old prior art
(Richards et al.\ 1990; Adamatzky's 1994 monograph)---the paper claims no
novelty for it. What nobody had measured is what that old fact \emph{does}
to a modern machine-learning benchmark. Answer: it quietly solves it.}

\subsection*{The benchmark (\S II): rules, targets, and what ``perfect'' means}

\textbf{What it does.} Fixes the observation protocol precisely (ring of
127 cells, 127 time steps, coin-flip random starting rows, diagrams that are
complete, noiseless, binary, and of known neighbourhood size), defines the
four damage statistics with their edge cases, and---the section's real
content---distinguishes three different ``ceilings'' a predictor can be held
to. Because each stored answer is one noisy measurement (estimated from
$K=20$ independent repeat measurements per rule): a predictor of the ideal
noise-free answer is capped by one ceiling (the ``ICC''), a predictor that
is itself a fresh noisy simulation is capped by a strictly lower one
(replicate agreement), and confusing the two misstates a method by more than
the margins being argued about. The $\rho$ score compresses this into one
number with fixed reference values ($\sqrt2$ and $1$) for every target.

\sowhat{Without this calibration, ``our model scores $0.88$'' is
uninterpretable---maybe $0.88$ is nearly perfect, maybe it is far off.
With it, every score in the paper can be read as a distance from what is
actually achievable. This reliability machinery is reusable in any study
whose ground truth comes from simulation.}

\textbf{Limits.} The reference values $\sqrt2$ and $1$ are exact only
\emph{on average} under an idealized noise model. The paper prices the
wobble: at most $0.14$ $\rho$-units from estimating noise with $K=20$
repeats, up to about $0.6$ from the finite rule panel---small against the
$1.6$--$34$-unit gaps the paper is about.

\subsection*{The three contestants (\S III)}

\textbf{What it does.} Describes the estimators. \emph{(i) Mechanistic:}
tabulate every pixel transition in the diagram, rebuild the rule table,
re-run the twin experiment on the rebuilt rule. Fully transparent, no
training, but it \emph{assumes} the observation model holds. \emph{(ii) Five
cheap statistics:} five simple hand-picked image summaries (density,
activity, compressibility, and two entropies) fed to gradient-boosted
decision trees. \emph{(iii) A small convolutional neural network} (CNN,
$31{,}524$ parameters) trained on many diagrams, with a first layer designed
to discourage rule-reading, five training repeats with different random
seeds, and---new in this version---weights chosen at the best epoch on an
inner slice of training data rather than simply the last epoch (a registered
protocol upgrade). The section also fixes the statistical rulebook in
advance: what counts as one method being ``superior'', ``equivalent'', or
the comparison being ``inconclusive''.

\sowhat{The contest is honest by construction: verdict rules, margins, and
test panels were committed to a timestamped repository before the
measurements they judge.}

\subsection*{Result 1 (\S IV A): rule reconstruction is as good as re-measuring}

\textbf{What it does.} Shows the mechanistic route recovers the exact rule
for $100\%$ of elementary and $97.5\%$ of radius-two held-out rules, and
proves that, once the rule is exactly recovered, its predictions and a fresh
independent measurement of the answer are statistically \emph{the same
thing}---so hitting the replicate ceiling is arithmetic, not luck. Scored
against a $16\times$ less noisy reference answer it rises to the top
ceiling, confirming its residual error was the comparison's noise, not its
own. The direct estimators do not rise when rescored the same way: their
shortfall is genuine. In headline numbers: median $R^2$ of $0.991$
(mechanistic) versus $0.880$ (CNN) and $0.845$ (trees) on the radius-two
panel---which sounds close until converted to noise units: $\rho$ of
$0.7$--$1.8$ versus $3$--$35$. Two deflating findings follow. The CNN's
seemingly excellent elementary-CA spreading-rate score ($R^2\approx0.98$)
is $\rho=20$: twenty times the noise floor. And on damage survival, most of
the direct estimators' apparent skill is \emph{mode assignment}: they
correctly sort rules into ``dies out'' versus ``spreads'' (the CNN gets
$0.94$ of elementary orbits and $1.00$ of radius-two rules right), but among
the genuinely uncertain middle rules their $R^2$ collapses (CNN $-1.05$ on
elementary CA---worse than guessing the group average). They are
classifiers that had been read as regressors. The section also discloses
that the radius-two test panel deliberately over-represents interesting
rules, and shows by decomposition that this works \emph{against} the paper's
own headline claim, not for it. The $2.5\%$ reconstruction failures are
examined and shown to be dynamically harmless.

\sowhat{This is the paper's central measurement: under clean observation,
system identification is simulation-limited---as good as re-running the
experiment---and the learned shortcuts are not remotely close, even where
their conventional scores look flattering.}

\subsection*{Result 2 (\S IV B): the two learned methods do not even beat each other}

\textbf{What it does.} Pre-registered paired comparisons show no stable
ranking between the CNN and the five-statistic trees: the CNN is ahead on
survival, the trees are superior on cone fill, the rest is inconclusive.
The mechanistic estimator meets the registered superiority bar against both
on most targets (differences up to $+0.39$ and $+0.53$) and loses to
neither on any target in either rule space.

\sowhat{If the expensive learned model and five cheap statistics trade
places depending on the target, neither has captured the underlying physics;
the method that read the rule has.}

\subsection*{Result 3 (\S IV C): what the network knows, and the memorisation test}

\textbf{What it does.} Two diagnostics. \emph{Stacking:} does the CNN carry
information the five statistics miss? Yes, on survival ($R^2$ increment
$+0.33$ over a simple base, $+0.09$ over the strong base)---but a
deployment-realistic combination cannot harvest it (increment $+0.001$, not
significant). \emph{Retrieval:} a lookup that simply returns the answer of
the most similar \emph{training} rule---no learning, no simulation---matches
the trained models ($0.864$ and $0.883$ versus the CNN's $0.880$; on
elementary CA the lookup inside the CNN's own internal representation
($0.883$) actually beats the CNN's own output ($0.847$)).

\sowhat{Much of the learned models' apparent skill is recognising which
training rule a diagram resembles---memorisation, not generalisation to
genuinely new rules. This cheap lookup test is portable to any benchmark
whose space of systems is densely sampled, and is one of the paper's most
exportable tools.}

\subsection*{Result 4 (\S IV D): the rule is readable everywhere, including inside the network}

\textbf{What it does.} Linear probes decode the exact rule identity from the
CNN's internal 64-number summary with $0.87$--$0.95$ accuracy (versus
near-zero chance), despite the anti-rule-reading first layer; individual
rule-table bits transfer to never-seen rules at $0.79$/$0.61$.

\sowhat{The information leak the paper warns about is not hypothetical: it
survives into the network trained to predict behaviour. The probes show the
information is \emph{present}; whether the network's predictions causally
\emph{use} it is explicitly left open.}

\subsection*{Result 5 (\S IV E--F): where reading fails, and what wins beyond}

\textbf{What it does.} Deliberately breaks each observation assumption and
maps the consequences on a grid (the ``identifiability frontier''). Exact
reconstruction is brittle exactly where you would expect: a deterministic
rule table has no concept of noise, so a single corrupted pixel transition
can flip a table entry---at $5\%$ pixel noise the mechanistic score falls
below zero. It tolerates about $25\%$ missing pixels and fails at extreme
starting densities. Then comes the repair race between eight registered
estimators plus two fairness controls: a \emph{probabilistic} rule reader
(the ``pseudo-posterior decoder''---keep a per-entry probability instead of
a hard bit, simulate several plausible tables, average) extends the working
range to $50\%$ masking ($R^2$ $0.85$ at $40\%$ masking versus the exact
inverter's $0.24$) and rescues extreme densities ($0.77$ versus $0.08$),
but only helps to about $3\%$ noise, and its confidence statements become
badly overconfident under noise (nominal $68\%$ intervals actually cover
$58\%$ at $2\%$ noise and $13\%$ at $20\%$). A learned rule-reading network
mostly fails---a memorable finding along the way: a table that is $73\%$
correct produces a \emph{badly wrong} dynamical system, because a rule
table is not the kind of object where ``mostly right'' helps. The fairness
controls matter most: retraining the same CNN \emph{on corrupted diagrams}
makes it the best method in the $3$--$15\%$ noise band ($R^2$
$0.51$--$0.62$)---so an earlier ``dead zone'' where nothing worked was an
artifact of testing clean-trained models---while an unconstrained network
$354\times$ larger changes no verdict, and a control that puts the
interesting rules back into training still leaves the CNN at $0.778$ versus
the mechanistic $0.992$. At $20\%$ noise nothing tested is reliable. All of
this is condensed into a practitioner guidance table: which estimator to use
in which observation regime, with its price.

\sowhat{This is the transferable product. The matched-regime verdict
(``just read the rule'') is specific to clean CA diagrams; the
\emph{frontier}---exact methods first, probabilistic reading under moderate
corruption, corruption-trained networks beyond, nothing past heavy
corruption---is the pattern the authors expect to generalize, and each axis
has a named analogue in real system identification (sensor noise, missing
sensors, insufficient excitation).}

\subsection*{Secondary analyses (\S IV, final subsections): cost, and honest by-products}

\textbf{What it does.} \emph{Cost:} the mechanistic estimator pays per
question ($283$\,ms) versus the CNN's $1$\,ms, but training means the CNN
only breaks even after ${\approx}4700$ questions (plausibly ${\approx}50$
against an optimized implementation)---and even a bargain-budget mechanistic
run ($17$\,ms) still beats both learned models on accuracy. A zero-cost
textbook approximation (the ``annealed'' mean-field theory) fails
completely ($R^2=-2.15$), for an identified reason: the targets measure
precisely the correlations that approximation throws away.
\emph{By-products:} a map of the radius-two rule space in damage coordinates
($7.8\%$ of rules meet a registered ``interesting damage'' criterion---
explicitly \emph{not} promoted to a validated census of complex rules), and
a sliding-window ``behaviour map'' for mixed-rule CA whose advantage is
statistically inconclusive. Both are reported at exactly the strength their
evaluation supports.

\sowhat{Speed is the only case left for the learned models in the matched
regime, and it is weaker than it looks. The by-products show the flip side
of the paper's discipline: results that did not pan out are reported as
such rather than dressed up.}

\subsection*{Discussion (\S V)}

\textbf{What it does.} States the claim at its correct scope: a
single-observation behaviour-prediction benchmark is solved by system
identification \emph{whenever} the observation pins down the generator
within a known model class, the inversion is tractable, and the defining
experiment can be re-run. Benchmark designers who intend to test
representation learning must break that chain (noise, partial observation,
unknown model class)---or accept that they are benchmarking identification.
The authors are careful about reach: they name no published benchmark this
invalidates, and they name their own most consequential omission---the
fully correct probabilistic decoder for noisy diagrams was not built, so
the noise-regime verdicts bound the estimators tested, not the whole
mechanistic family. The identified next experiment: stochastic CA
(Domany--Kinzel), where the rule itself is probabilistic, identification
becomes genuine estimation, and the paper's near-identity argument should
finally break down.

\sowhat{The discussion converts one clean case study into a design rule for
a whole literature, while drawing the boundary of its own evidence
unusually sharply.}

\subsection*{Appendices A--E}

\emph{A: Reproducibility and registration.} The full ledger of which
analyses were registered before measurement, which were added afterwards,
and which are exploratory; plus a correctness audit of the central identity
and one disclosed deviation from the original registration. \emph{B:
Methods supplement.} Exact data, splits, network architecture, evaluation
protocol, and the record of the checkpoint-selection switch (the registered
rule fired: radius-two cone fill rose $0.337\to0.406$ and seed-to-seed
spread collapsed $0.109\to0.036$). \emph{C: Decoder algorithms.} Pseudocode
for every rule-reading variant, with a worked example showing why the
probabilistic decoder becomes overconfident. \emph{D: The damage-signature
landscape.} The rule-space map with its central caveat spelled out: the
striking-looking region is substantially a coordinate effect, a registered
threshold region rather than a discovered cluster; an intended independent
detector performed at chance and validates nothing. \emph{E: Prior work
under its assumptions.} A table of thirty-five years of CA identification
literature, so the boundary of the novelty claim is auditable.

\section*{5. The limits, stated plainly}

\begin{itemize}
\item \textbf{Everything is protocol-scoped.} All statistics are defined for
this exact setup: ring of 127 cells, 127 steps, coin-flip starting rows.
This is deliberate (the asymptotic versions are formally undecidable), but
it means no claim is about ``the rule'' in the abstract.
\item \textbf{The neighbourhood size is assumed known} in the main analysis;
inferring it costs about $0.12$ in median $R^2$.
\item \textbf{The noise-regime verdicts bound the tested estimators only.}
The correctly specified probabilistic decoder for noisy diagrams was not
built; the authors flag this as their most consequential omission.
\item \textbf{Neural claims are scoped to this architecture and budget.} No
architecture search was run; the paper explicitly does not conclude that
deep networks cannot close the gap.
\item \textbf{The radius-two test panel is enriched} with interesting rules
selected using the same data that supplies the evaluation labels---
disclosed, and shown by two controls to be conservative for the headline
claim, but a design wart nonetheless.
\item \textbf{Deterministic, binary, one-dimensional CA only.} The
stochastic case is argued to be the thesis's breaking point---and has not
been run.
\item \textbf{The frontier rests on a fixed 80-rule panel}; only the noise
axis was replicated on a second panel (the qualitative structure held, and
the crossover moved---it is a band, not a line); corruptions were tested one
axis at a time.
\item \textbf{Timings are one machine, one implementation.}
\end{itemize}

\section*{6. Why the results are valuable anyway}

\textbf{A calibrated negative result.} Most of the paper's value is in what
it stops people from doing: publishing learned single-observation predictors
of perturbation response without building the identification baseline. The
gap is not asserted but measured, in units that make it unmistakable.

\textbf{Reusable measurement machinery.} The three-ceiling reliability
analysis and the $\rho$ score apply to any benchmark whose ground truth
comes from noisy simulation---which is most of simulation-based machine
learning.

\textbf{Two cheap, portable diagnostics.} The mode-assignment check (is the
regressor secretly a classifier?) and the nearest-training-example lookup
(is the model secretly memorising?) can be run on any benchmark in an
afternoon.

\textbf{A practical map.} The frontier and its guidance table tell a
practitioner which estimator to trust for a given observation quality, with
measured evidence and explicit refusal to recommend anything where nothing
worked.

\textbf{A methodological standard.} Decision rules were frozen before
measurements, every quoted number is re-derived from released artifacts by
an audit script, negative and inconclusive results are reported at face
value, and each claim is labelled registered, post-hoc, or exploratory. The
paper practises the benchmark hygiene it preaches.

\vspace{6pt}
\noindent\rule{\textwidth}{0.4pt}
\vspace{2pt}

\noindent\small This summary accompanies the repository copy of the
manuscript. It simplifies wording but adds no claims; where it paraphrases,
the manuscript's own hedges apply. Source of record:
\texttt{manuscript/main.tex} at commit \texttt{8bf5b4a}.

\end{document}
